\documentclass[10pt,a4paper]{amsart}
\usepackage{amsmath,amssymb,amsthm,graphicx}
% Retained minimal preamble for the Stacks Project chapter more-algebra.tex (More on Algebra).
% The chapter file pinned-original-source.tex carries only \input{preamble} on its first line, so no
% preamble text lives inside the chapter itself.  The Stacks Project uses one single shared preamble for
% every chapter; the snapshot stages it identically as sources/stacks-chow/preamble.tex and
% sources/stacks-etale-cohomology/preamble.tex.  The block below is copied mechanically, line for line,
% from sources/stacks-chow/preamble.tex lines 14-18 (the Xy-pic setup) and 154-217 (the native theorem
% declarations together with the mathematical macros).  Nothing was added, removed or redefined; the
% wrapper class and the amsmath/amssymb/amsthm/graphicx base packages are supplied by the assembler.
% Chapter: more-algebra.tex -- More on Algebra.
% For commutative diagrams we use Xy-pic
\usepackage[all]{xy}

% We use 2cell for 2-commutative diagrams.
\xyoption{2cell}

% Theorem environments.
%
\theoremstyle{plain}
\newtheorem{theorem}[subsection]{Theorem}
\newtheorem{proposition}[subsection]{Proposition}
\newtheorem{lemma}[subsection]{Lemma}

\theoremstyle{definition}
\newtheorem{definition}[subsection]{Definition}
\newtheorem{example}[subsection]{Example}
\newtheorem{exercise}[subsection]{Exercise}
\newtheorem{situation}[subsection]{Situation}

\theoremstyle{remark}
\newtheorem{remark}[subsection]{Remark}
\newtheorem{remarks}[subsection]{Remarks}

\numberwithin{equation}{subsection}

% Macros
%
\def\lim{\mathop{\mathrm{lim}}\nolimits}
\def\colim{\mathop{\mathrm{colim}}\nolimits}
\def\Spec{\mathop{\mathrm{Spec}}}
\def\Hom{\mathop{\mathrm{Hom}}\nolimits}
\def\Ext{\mathop{\mathrm{Ext}}\nolimits}
\def\SheafHom{\mathop{\mathcal{H}\!\mathit{om}}\nolimits}
\def\SheafExt{\mathop{\mathcal{E}\!\mathit{xt}}\nolimits}
\def\Sch{\mathit{Sch}}
\def\Mor{\mathop{\mathrm{Mor}}\nolimits}
\def\Ob{\mathop{\mathrm{Ob}}\nolimits}
\def\Sh{\mathop{\mathit{Sh}}\nolimits}
\def\NL{\mathop{N\!L}\nolimits}
\def\CH{\mathop{\mathrm{CH}}\nolimits}
\def\proetale{{pro\text{-}\acute{e}tale}}
\def\etale{{\acute{e}tale}}
\def\QCoh{\mathit{QCoh}}
\def\Ker{\mathop{\mathrm{Ker}}}
\def\Im{\mathop{\mathrm{Im}}}
\def\Coker{\mathop{\mathrm{Coker}}}
\def\Coim{\mathop{\mathrm{Coim}}}

% Boxtimes
%
\DeclareMathSymbol{\boxtimes}{\mathbin}{AMSa}{"02}

%
% Macros for moduli stacks/spaces
%
\def\QCohstack{\mathcal{QC}\!\mathit{oh}}
\def\Cohstack{\mathcal{C}\!\mathit{oh}}
\def\Spacesstack{\mathcal{S}\!\mathit{paces}}
\def\Quotfunctor{\mathrm{Quot}}
\def\Hilbfunctor{\mathrm{Hilb}}
\def\Curvesstack{\mathcal{C}\!\mathit{urves}}
\def\Polarizedstack{\mathcal{P}\!\mathit{olarized}}
\def\Complexesstack{\mathcal{C}\!\mathit{omplexes}}
% \Pic is the operator that assigns to X its picard group, usage \Pic(X)
% \Picardstack_{X/B} denotes the Picard stack of X over B
% \Picardfunctor_{X/B} denotes the Picard functor of X over B
\def\Pic{\mathop{\mathrm{Pic}}\nolimits}
\def\Picardstack{\mathcal{P}\!\mathit{ic}}
\def\Picardfunctor{\mathrm{Pic}}
\def\Deformationcategory{\mathcal{D}\!\mathit{ef}}
\usepackage{cleveref}
\crefname{originallemma}{original lemma}{original lemmas}
\Crefname{originallemma}{Original lemma}{Original lemmas}
\makeatletter
\newlabel{algebra-lemma-Noetherian-complete-local-Nagata}{{032W}{source}{}{sourcepointer-1}{}}
\newlabel{algebra-lemma-Noetherian-complete-local-Nagata@cref}{{[originallemma][0][]032W}{[1][0][]source}{}{}{}}
\makeatother
\crefname{originallemma}{original lemma}{original lemmas}
\Crefname{originallemma}{Original lemma}{Original lemmas}
\makeatletter
\newlabel{algebra-section-completion}{{00M9}{source}{}{sourcepointer-2}{}}
\newlabel{algebra-section-completion@cref}{{[originallemma][0][]00M9}{[1][0][]source}{}{}{}}
\makeatother
\crefname{originallemma}{original lemma}{original lemmas}
\Crefname{originallemma}{Original lemma}{Original lemmas}
\makeatletter
\newlabel{fields-lemma-normal-case}{{030M}{source}{}{sourcepointer-3}{}}
\newlabel{fields-lemma-normal-case@cref}{{[originallemma][0][]030M}{[1][0][]source}{}{}{}}
\makeatother
\crefname{originallemma}{original lemma}{original lemmas}
\Crefname{originallemma}{Original lemma}{Original lemmas}
\makeatletter
\newlabel{fields-lemma-normal-closure}{{09DT}{source}{}{sourcepointer-4}{}}
\newlabel{fields-lemma-normal-closure@cref}{{[originallemma][0][]09DT}{[1][0][]source}{}{}{}}
\makeatother
\crefname{originallemma}{original lemma}{original lemmas}
\Crefname{originallemma}{Original lemma}{Original lemmas}
\makeatletter
\newlabel{lemma-abhyankar}{{34037}{source}{}{sourcepointer-5}{}}
\newlabel{lemma-abhyankar@cref}{{[originallemma][0][]34037}{[1][0][]source}{}{}{}}
\makeatother
\crefname{originallemma}{original lemma}{original lemmas}
\Crefname{originallemma}{Original lemma}{Original lemmas}
\makeatletter
\newlabel{lemma-characteristic-p-case}{{34818}{source}{}{sourcepointer-6}{}}
\newlabel{lemma-characteristic-p-case@cref}{{[originallemma][0][]34818}{[1][0][]source}{}{}{}}
\makeatother
\crefname{originallemma}{original lemma}{original lemmas}
\Crefname{originallemma}{Original lemma}{Original lemmas}
\makeatletter
\newlabel{lemma-galois-inertia}{{33251}{source}{}{sourcepointer-7}{}}
\newlabel{lemma-galois-inertia@cref}{{[originallemma][0][]33251}{[1][0][]source}{}{}{}}
\makeatother
\crefname{originallemma}{original lemma}{original lemmas}
\Crefname{originallemma}{Original lemma}{Original lemmas}
\makeatletter
\newlabel{lemma-galois-relative}{{34609}{source}{}{sourcepointer-8}{}}
\newlabel{lemma-galois-relative@cref}{{[originallemma][0][]34609}{[1][0][]source}{}{}{}}
\makeatother
\crefname{originallemma}{original lemma}{original lemmas}
\Crefname{originallemma}{Original lemma}{Original lemmas}
\makeatletter
\newlabel{lemma-inertial-invariants-unramified}{{33475}{source}{}{sourcepointer-9}{}}
\newlabel{lemma-inertial-invariants-unramified@cref}{{[originallemma][0][]33475}{[1][0][]source}{}{}{}}
\makeatother
\crefname{originallemma}{original lemma}{original lemmas}
\Crefname{originallemma}{Original lemma}{Original lemmas}
\makeatletter
\newlabel{lemma-lowering-the-level}{{35165}{source}{}{sourcepointer-10}{}}
\newlabel{lemma-lowering-the-level@cref}{{[originallemma][0][]35165}{[1][0][]source}{}{}{}}
\makeatother
\crefname{originallemma}{original lemma}{original lemmas}
\Crefname{originallemma}{Original lemma}{Original lemmas}
\makeatletter
\newlabel{lemma-make-finite-level}{{35088}{source}{}{sourcepointer-11}{}}
\newlabel{lemma-make-finite-level@cref}{{[originallemma][0][]35088}{[1][0][]source}{}{}{}}
\makeatother
\crefname{originallemma}{original lemma}{original lemmas}
\Crefname{originallemma}{Original lemma}{Original lemmas}
\makeatletter
\newlabel{lemma-purely-inseparable-case}{{34676}{source}{}{sourcepointer-12}{}}
\newlabel{lemma-purely-inseparable-case@cref}{{[originallemma][0][]34676}{[1][0][]source}{}{}{}}
\makeatother
\crefname{originallemma}{original lemma}{original lemmas}
\Crefname{originallemma}{Original lemma}{Original lemmas}
\makeatletter
\newlabel{lemma-solutions-go-down}{{34517}{source}{}{sourcepointer-13}{}}
\newlabel{lemma-solutions-go-down@cref}{{[originallemma][0][]34517}{[1][0][]source}{}{}{}}
\makeatother
\crefname{originallemma}{original lemma}{original lemmas}
\Crefname{originallemma}{Original lemma}{Original lemmas}
\makeatletter
\newlabel{lemma-weakly-unramified-goes-up-along-totally-ramified}{{34490}{source}{}{sourcepointer-14}{}}
\newlabel{lemma-weakly-unramified-goes-up-along-totally-ramified@cref}{{[originallemma][0][]34490}{[1][0][]source}{}{}{}}
\makeatother
\crefname{originalremark}{original remark}{original remarks}
\Crefname{originalremark}{Original remark}{Original remarks}
\makeatletter
\newlabel{remark-construction}{{33882}{source}{}{sourcepointer-15}{}}
\newlabel{remark-construction@cref}{{[originalremark][0][]33882}{[1][0][]source}{}{}{}}
\makeatother
\title{For an algebraically closed residue field in characteristic p and complete weakly unramified rings, weak solutions exist}
\author{The Stacks Project Authors}
\date{stacks-project:more-algebra.tex:local-snapshot-20261012; GNU FDL 1.2 or later}
\begin{document}
\maketitle
\setcounter{subsection}{1}
\noindent\textit{Editorial note.} All delivered mathematical text is the author's own, copied from the pinned Stacks Project chapter file more-algebra.tex; no statement or proof step was rewritten, reordered or removed.\par\medskip
\noindent\textit{Editorial note.} The chapter file pinned-original-source.tex carries only \textbackslash{}input\{preamble\} on its first line, so no preamble text lives inside the chapter. The retained file sources/stacks-more-algebra/retained-preamble.tex is a mechanical, line-for-line copy of the project's single shared preamble as staged in the snapshot; nothing was added, removed or redefined and no package was added.\par\medskip
\noindent\textit{Editorial note.} The source URL is the source-level Stacks Project URL, as required when more than one card is produced from one source; the per-claim tag is carried by claim\_key.\par\medskip
\noindent\textit{Editorial note.} The counted proof span is the author's complete proof environment at source lines 35385-35511. Its references are recorded as non-load-bearing navigation pointers; both the cross-chapter ones (resolving to the project's tag list) and the same-chapter ones (resolving to their \textbackslash{}label lines in the pinned source) concern standard structural facts whose applications stay inside the counted proof span.\par\medskip
\noindent\textit{Allowed standard prerequisites (editorial).}\begin{itemize}
\item Nagata rings and finiteness of integral closures in finite extensions. invoked at proof lines 35387-35388. Is what makes the integral closure of B in a finite extension finite over B, the first step of the proof.
\item Valuation rings as maximal with respect to domination, and completeness of a finitely generated module over a complete ring. used at proof lines 35389-35395. Identifies C with the integral closure of B in M and shows the auxiliary rings are complete.
\item Inertia and ramification structure theory of extensions of discrete valuation rings. developed earlier in the same section and used throughout the rest of the counted proof span. Supplies the reduction to the totally ramified case in which the characteristic p machinery applies.
\end{itemize}\medskip
\ifdefined\hypertarget\hypertarget{sourcepointer-1}{}\fi
\noindent\textit{Source reference.} Original lemma 032W: focused-seeds/AGP/sources/stacks-snapshot-20261011/tags.txt, source lines 4035--4035. Cross-chapter navigation pointer into the Stacks Project; the standard fact it supplies is applied inside the counted proof span and the argument resting on it stays there.\par\medskip
\ifdefined\hypertarget\hypertarget{sourcepointer-2}{}\fi
\noindent\textit{Source reference.} Original lemma 00M9: focused-seeds/AGP/sources/stacks-snapshot-20261011/tags.txt, source lines 815--815. Cross-chapter navigation pointer into the Stacks Project; the standard fact it supplies is applied inside the counted proof span and the argument resting on it stays there.\par\medskip
\ifdefined\hypertarget\hypertarget{sourcepointer-3}{}\fi
\noindent\textit{Source reference.} Original lemma 030M: focused-seeds/AGP/sources/stacks-snapshot-20261011/tags.txt, source lines 3953--3953. Cross-chapter navigation pointer into the Stacks Project; the standard fact it supplies is applied inside the counted proof span and the argument resting on it stays there.\par\medskip
\ifdefined\hypertarget\hypertarget{sourcepointer-4}{}\fi
\noindent\textit{Source reference.} Original lemma 09DT: focused-seeds/AGP/sources/stacks-snapshot-20261011/tags.txt, source lines 11876--11876. Cross-chapter navigation pointer into the Stacks Project; the standard fact it supplies is applied inside the counted proof span and the argument resting on it stays there.\par\medskip
\ifdefined\hypertarget\hypertarget{sourcepointer-5}{}\fi
\noindent\textit{Source reference.} Original lemma 34037: sources/stacks-more-algebra/pinned-original-source.tex, source lines 34037--34037. Same-chapter navigation pointer; the statement it names is standard for the section and the argument that consumes it stays inside the counted proof span.\par\medskip
\ifdefined\hypertarget\hypertarget{sourcepointer-6}{}\fi
\noindent\textit{Source reference.} Original lemma 34818: sources/stacks-more-algebra/pinned-original-source.tex, source lines 34818--34818. Same-chapter navigation pointer; the statement it names is standard for the section and the argument that consumes it stays inside the counted proof span.\par\medskip
\ifdefined\hypertarget\hypertarget{sourcepointer-7}{}\fi
\noindent\textit{Source reference.} Original lemma 33251: sources/stacks-more-algebra/pinned-original-source.tex, source lines 33251--33251. Same-chapter navigation pointer; the statement it names is standard for the section and the argument that consumes it stays inside the counted proof span.\par\medskip
\ifdefined\hypertarget\hypertarget{sourcepointer-8}{}\fi
\noindent\textit{Source reference.} Original lemma 34609: sources/stacks-more-algebra/pinned-original-source.tex, source lines 34609--34609. Same-chapter navigation pointer; the statement it names is standard for the section and the argument that consumes it stays inside the counted proof span.\par\medskip
\ifdefined\hypertarget\hypertarget{sourcepointer-9}{}\fi
\noindent\textit{Source reference.} Original lemma 33475: sources/stacks-more-algebra/pinned-original-source.tex, source lines 33475--33475. Same-chapter navigation pointer; the statement it names is standard for the section and the argument that consumes it stays inside the counted proof span.\par\medskip
\ifdefined\hypertarget\hypertarget{sourcepointer-10}{}\fi
\noindent\textit{Source reference.} Original lemma 35165: sources/stacks-more-algebra/pinned-original-source.tex, source lines 35165--35165. Same-chapter navigation pointer; the statement it names is standard for the section and the argument that consumes it stays inside the counted proof span.\par\medskip
\ifdefined\hypertarget\hypertarget{sourcepointer-11}{}\fi
\noindent\textit{Source reference.} Original lemma 35088: sources/stacks-more-algebra/pinned-original-source.tex, source lines 35088--35088. Same-chapter navigation pointer; the statement it names is standard for the section and the argument that consumes it stays inside the counted proof span.\par\medskip
\ifdefined\hypertarget\hypertarget{sourcepointer-12}{}\fi
\noindent\textit{Source reference.} Original lemma 34676: sources/stacks-more-algebra/pinned-original-source.tex, source lines 34676--34676. Same-chapter navigation pointer; the statement it names is standard for the section and the argument that consumes it stays inside the counted proof span.\par\medskip
\ifdefined\hypertarget\hypertarget{sourcepointer-13}{}\fi
\noindent\textit{Source reference.} Original lemma 34517: sources/stacks-more-algebra/pinned-original-source.tex, source lines 34517--34517. Same-chapter navigation pointer; the statement it names is standard for the section and the argument that consumes it stays inside the counted proof span.\par\medskip
\ifdefined\hypertarget\hypertarget{sourcepointer-14}{}\fi
\noindent\textit{Source reference.} Original lemma 34490: sources/stacks-more-algebra/pinned-original-source.tex, source lines 34490--34490. Same-chapter navigation pointer; the statement it names is standard for the section and the argument that consumes it stays inside the counted proof span.\par\medskip
\ifdefined\hypertarget\hypertarget{sourcepointer-15}{}\fi
\noindent\textit{Source reference.} Original remark 33882: sources/stacks-more-algebra/pinned-original-source.tex, source lines 33882--33882. Same-chapter navigation pointer; the statement it names is standard for the section and the argument that consumes it stays inside the counted proof span.\par\medskip
\begin{lemma}
\label{lemma-special-case}
Let $A \subset B \subset C$ be extensions of discrete valuation rings
with fraction fields $K \subset L \subset M$. Assume
\begin{enumerate}
\item the residue field $k$ of $A$ is algebraically closed of
characteristic $p > 0$,
\item $A$ and $B$ are complete,
\item $A \to B$ is weakly unramified,
\item $M$ is a finite extension of $L$,
\item $k = \bigcap\nolimits_{n \geq 1} \kappa_B^{p^n}$
\end{enumerate}
Then there exists a finite extension $K_1/K$ which
is a weak solution for $A \to C$.
\end{lemma}

\begin{proof}
Let $M'$ be any finite extension of $L$ and consider the integral closure
$C'$ of $B$ in $M'$. Then $C'$ is finite over $B$ as $B$ is Nagata by
Algebra, Lemma \ref{algebra-lemma-Noetherian-complete-local-Nagata}.
Moreover, $C'$ is a discrete valuation ring, see discussion in
Remark \ref{remark-construction}. Moreover $C'$ is complete as a
$B$-module, hence complete as a discrete valuation ring, see
Algebra, Section \ref{algebra-section-completion}.
It follows in particular that $C$ is the integral
closure of $B$ in $M$ (by definition of valuation rings as maximal
for the relation of domination).

\medskip\noindent
Let $M \subset M'$ be a finite extension and let $C' \subset M'$
be the integral closure of $B$ as above. By
Lemma \ref{lemma-solutions-go-down}
it suffices to prove the result for $A \to B \to C'$.
Hence we may assume that $M/L$ is normal, see
Fields, Lemma \ref{fields-lemma-normal-closure}.

\medskip\noindent
If $M / L$ is normal, we can find a chain of finite extensions
$$
L = L^0 \subset L^1 \subset L^2 \subset \ldots \subset L^r = M
$$
such that each extension $L^{j + 1}/L^j$ is either:
\begin{enumerate}
\item[(a)] purely inseparable of degree $p$,
\item[(b)] totally ramified with respect to $B^j$ and Galois of degree $p$,
\item[(c)] totally ramified with respect to $B^j$ and Galois cyclic of
order prime to $p$,
\item[(d)] Galois and unramified with respect to $B^j$.
\end{enumerate}
Here $B^j$ is the integral closure of $B$ in $L^j$.
Namely, since $M/L$ is normal we can write it as a compositum of
a Galois extension and a purely inseparable extension
(Fields, Lemma \ref{fields-lemma-normal-case}).
For the purely inseparable extension the existence of the filtration
is clear. In the Galois case, note that $G$ is ``the'' decomposition group
and let $I \subset G$ be the inertia group. Then on the one hand
$I$ is solvable by Lemma \ref{lemma-galois-inertia} and on the other
hand the extension $M^I/L$ is unramified with respect to $B$ by
Lemma \ref{lemma-inertial-invariants-unramified}.
This proves we have a filtration as stated.

\medskip\noindent
We are going to argue by induction on the integer $r$. Suppose that we
can find a finite extension $K_1/K$ which is a weak solution
for $A \to B^1$ where $B^1$ is the integral closure of $B$ in $L^1$.
Let $K'_1$ be the normal closure of $K_1/K$
(Fields, Lemma \ref{fields-lemma-normal-closure}).
Since $A$ is complete and the residue field of $A$ is algebraically closed
we see that $K'_1/K_1$ is separable and totally ramified with
respect to $A_1$ (some details omitted).
Hence $K'_1/K$ is a weak solution for $A \to B^1$ as well by
Lemma \ref{lemma-weakly-unramified-goes-up-along-totally-ramified}.
In other words, we may and do assume that $K_1$ is a normal extension of $K$.
Having done so we consider the sequence
$$
L^0_1 = (L^0 \otimes_K K_1)_{red} \subset
L^1_1 = (L^1 \otimes_K K_1)_{red} \subset \ldots \subset
L^r_1 = (L^r \otimes_K K_1)_{red}
$$
and the corresponding integral closures $B^i_1$. Note that $C_1 = B^r_1$
is a product of discrete valuation rings which are transitively permuted
by $G = \text{Aut}(K_1/K)$ by Lemma \ref{lemma-galois-relative}.
In particular all the extensions of discrete valuation rings
$A_1 \to (C_1)_\mathfrak m$ are isomorphic and a weak solution for one
will be a weak solution for all of them. We can apply the induction
hypothesis to the sequence
$$
A_1 \to (B^1_1)_{B^1_1 \cap \mathfrak m} \to
(B^2_1)_{B^2_1 \cap \mathfrak m} \to
\ldots \to
(B^r_1)_{B^r_1 \cap \mathfrak m} =
(C_1)_\mathfrak m
$$
to get a weak solution $K_2/K_1$ for $A_1 \to (C_1)_\mathfrak m$.
The extension $K_2/K$ will then be a weak solution for $A \to C$
by what we said before. Note that the induction hypothesis applies:
the ring map $A_1 \to (B^1_1)_{B^1_1 \cap \mathfrak m}$
is weakly unramified by our choice of $K_1$
and the sequence of fraction field extensions
each still have one of the properties (a), (b), (c), or (d)
listed above. Moreover, observe that for any finite extension 
$\kappa_B \subset \kappa$ we still have $k = \bigcap \kappa^{p^n}$.

\medskip\noindent
Thus everything boils down to finding a weak solution for $A \subset C$
when the field extension $M/L$ satisfies one of the properties
(a), (b), (c), or (d).

\medskip\noindent
Case (d). This case is trivial as here $B \to C$ is unramified already.

\medskip\noindent
Case (c). Say $M/L$ is cyclic of order $n$ prime to $p$. Because
$M/L$ is totally ramified with respect to $B$ we see that the ramification
index of $B \subset C$ is $n$ and hence the ramification index of $A \subset C$
is $n$ as well. Choose a uniformizer $\pi \in A$ and set
$K_1 = K[\pi^{1/n}]$. Then $K_1/K$ is a solution for $A \subset C$
by Abhyankar's lemma (Lemma \ref{lemma-abhyankar}).

\medskip\noindent
Case (b). We divide this case into the mixed characteristic case and the
equicharacteristic case. In the equicharacteristic case this is
Lemma \ref{lemma-characteristic-p-case}. In the mixed characteristic
case, we first replace $K$ by a finite extension to get to the
situation where $M/L$ is a degree $p$ extension of finite level using
Lemma \ref{lemma-make-finite-level}.
Then the level is a rational number $l \in [0, p)$, see discussion
preceding Lemma \ref{lemma-lowering-the-level}. If the level is $0$,
then $B \to C$ is weakly unramified and we're done. If not, then we
can replacing the field $K$ by a finite extension to obtain a new
situation with level $l' \leq \max(0, l - 1, 2l - p)$ by
Lemma \ref{lemma-lowering-the-level}.
If $l = p - \epsilon$ for $\epsilon < 1$ then we see that
$l' \leq p - 2\epsilon$. Hence after a finite number of replacements
we obtain a case with level $\leq p - 1$. Then after at most $p - 1$
more such replacements we reach the situation where the level is zero.

\medskip\noindent
Case (a) is Lemma \ref{lemma-purely-inseparable-case}. This is the only case
where we possibly need a purely inseparable extension of $K$, namely, in
case (2) of the statement of the lemma we win by adjoining a $p$th power
of the element $\pi$. This finishes the proof of the lemma.
\end{proof}

\end{document}
